\documentclass[12px]{article}

\title{Lezione 20 Metodi}
\date{2026-05-29}
\author{Federico De Sisti}

\input{../../setup.tex}

\begin{document}
	\maketitle
	\newpage
	\subsection{Zero-Stabilità dei metodi di Adams}
	Zero-Stabilità $ \Leftrightarrow$ condizione delle radici su $\rho(r)$\\
	 \[
		 u_{n+1} = u_n + h \sum^{p}_{j=-1}b_j f_{n-j} \ \ \ \rho(r) = r^{p+1} - r^p = r^p(r-1)
	.\] 
	Le radici sono: $r=1$ semplice,  $r=0$ con molteplitictà $p$ $ \Rightarrow  $, il criterio della radice è rispettato\\
	\textbf{Osservazione}\\
	Per l'assoluta stabilità bisogna calcolare le radici di  $P(r) = \rho(r) - z\sigma(r)$ e verificare che  $|r_j(z)|<1$\\
	In generale la forma del polinomio $P(r)$ non fornisce le radici esplicitamente.\\
	\begin{defi}[Regione di stabilità]
		Chiamo $A$ la regione di stabilità
		 \[
			 A = \{z\in \C,Re(z) < 0 \ \ t.c.\ \ |r_j(z)| < 1, \ j = 0,\ldots, p\}
		.\] 
		definisco poi
		\[
			R_j = \{z\in \C, Re(z) < 0 \ \ t.c\ \ |r_j(z)| < 1\}
		.\] \\
		Allora
		 \[
			 A = \bigcap^p_{j=0}R_j
		.\] 
	\end{defi}
	Vedremo un modo più facile per visualizzare $A$
	\subsection{Metodi BDF (Backward Differentiation Formule)}
	Sono particolari metodi LMM sempre impliciti della forma
	 \[
		 u_{n+1} = \sum^{p}_{j=0}a_ju_{n-j} + h b_{-1}f_{n+1}
	.\] 
	Introdotti per trattare i cosiddetti problemi SIFF (problemi che presentano autovalori $\lambda$ con moduli molto diversi fra loro)\\
	\textbf{Osservazione}\\[5px]
	$\displaystyle \frac{u_{n+1} - \sum^{p}_{j=0}a_ju_{n-j}}{hb_{-1}} = f_{n+1} = f(t_{n+1},u_{n+1}) \xrightarrow{h \rightarrow 0} f(t_{n+1},y_{n+1}) = \dot y(t_{n+1})$\\
	Quindi i coefficienti devono essere scelti per approssimare $\dt y(t_{n+1})$ tramite una interpolazione di  $y$ nei nodi $t_{n-p},\ldots, t_n, t_{n+1}$ che sono  $p+2$\\
	Questo si fa con un polinomio di grado $p+1$\\
	Una volta calcolato  $\Pi_{p+1}[y][t_{n-p},\ldots,t_{n+1}](t)$ si calcola la sua derivata prima  $\Pi_{p+1}'$ che è un approssimazione di  $\dot y(t)$ e la si calcola con $t = t_{n+1}$ ottenendo lo schema BDF:  $\underbrace{\Pi'_{p+1}[u_{n-p},\ldots,u_{n+1}] = f(t_{n+1},u_{n+1})}_{\text{impongo}}$\\
	\textbf{Esempio} $ p = 0$\\
\[\Pi_1[u_n,u_{n+1}](t) = u_n + (t-t_n)\frac{u_{n+1}-u_n}{t_{n+1}-t_n}\]
\[
	\Pi_1'[u_n,u_{n+1}](t) = \frac{ u_{n+1} - u_n}{t_{n+1}-t_n} = \frac {u_{n+1}-u_n}h
.\] 
\[
	\leadsto \Pi'_1[u_{n},u_{n+1}](t_{n+1}) = \frac{u_{n+1} - u_n}h \Rightarrow \frac{u_{n+1}-u_n}{h} = f(t_{n+1},u_{n+1}) \ \ b_{-1} =1
.\] 
$ \Rightarrow BDF1 \equiv BE$ \\
\textbf{BDF2}: $u_{n-1},u_n,u_{n+1}$ Costruisco il polinomio interpolatorio  $\Pi_2(t)$\\
Useremo la forma di Newton $\Pi_2(t)$ (Differenze divise)\\


\[
	\begin{centered}
		\hspace{-40px}
\begin{tikzpicture}
  % We apply \small to shrink the text size slightly if needed,
  % but adjusting column sep is usually enough.
  \matrix (m) [
    matrix of math nodes,
    row sep=1em,        % Reduced vertical space
    column sep=0.5em,   % Drastically reduced horizontal space
    nodes={anchor=center}
  ] {
    % Riga 1
    & \textcolor{red}{\text{Diff div ORD 0}} & \textcolor{red}{\text{ORD 1}} & \textcolor{red}{\text{ORD 2}} \\
    % Riga 2
    t_{n-1} : & u_{n-1} & & \\
    % Riga 3
    & & \displaystyle u[t_{n-1}, t_n] = \frac{u_n - u_{n-1}}{h} & \\
    % Riga 4
    t_n : & u_n & & \displaystyle u[t_{n-1}, t_n, t_{n+1}] = \frac{u[t_n, t_{n+1}] - u[t_{n-1}, t_n]}{2h} \\
    % Riga 5
    & & \displaystyle u[t_n, t_{n+1}] = \frac{u_{n+1} - u_n}{h} & \\
    % Riga 6
    t_{n+1} : & u_{n+1} & & \\
  };

  % Arrows
  \begin{scope}[every path/.style={->, shorten >=3pt, shorten <=3pt}]
    % Da ORD 0 a ORD 1
    \draw (m-2-2.east) -- (m-3-3.west);
    \draw (m-4-2.east) -- (m-3-3.west);

    \draw (m-4-2.east) -- (m-5-3.west);
    \draw (m-6-2.east) -- (m-5-3.west);

    % Da ORD 1 a ORD 2
    \draw (m-3-3.east) -- (m-4-4.west);
    \draw (m-5-3.east) -- (m-4-4.west);
  \end{scope}

\end{tikzpicture}
	\end{centered}
\]
\[
	\Rightarrow  u[t_{n-1},t_n,t_{n+1}] = \frac{ \frac{u_{n+1}-u_n}{h} - \frac{u_n - u_{n-1}}{h}}{2h} = \frac{u_{n+1} - 2u_n + u_{n-1}}{2h^2}
.\] 
\[
	\begin{split}
		\Rightarrow \Pi_2(t) &= u_{n+1} + u[t_{n-1},t_n](t-t_{n-1}) + u[t_{n-1},t_n,t_{n+1}](t-t_{n-1})(t-t_n)\\
				     &= u_{n+1} + \frac{u_n - u_{n-1}}{h}(t-t_{n-1}) + \frac{u_{n+1}-2u_n + u_{n-1}}{2h^2}(t^2-t_{n-1}t - t_nt + t_nt_{n-1})\\
				     &= u_{n+1} + \frac{u_n - u_{n-1}}{h}(t-t_{n-1}) + \frac{u_{n+1}-2u_n + u_{n-1}}{2h^2}(t-t_{n-1})(t-t_n)
	\end{split}
.\] 
\[
	\Pi_2'(t) = \frac{u_n - u_{n-1}}{h} + \frac{u_{n+1} - 2u_n + u_{n-1}}{2h^2}(t-t_n + t-t_{n-1})
\] 
\[
	\Pi'_2(t_{n+1}) = \frac{u_n - u_{n-1}}{h} + \frac{u_{n+1} - 2u_n + u_{n-1}}{2h^{\cancel 2}}( \overbrace{t_{n-1}- t_n}^{\cancel h} +\overbrace{t_{n+1}-t_{n-1}}^{2\cancel h} ) \underbrace{=}_{\text{impongo}} f(t_{n+1}, u_{n+1})
.\] 
Quindi
\[
	\frac{u_n-u_{n-1}}{h} + \frac{u_{n+1} - 2u_n + u_{n-1}}{2h}\cdot 3 = f(t_{n+1},u_{n+1})
.\] 
\[
2(u_n - u_{n-1}) + 3(u_{n+1} - 2u_n + u_{n+1}) = 2h f(t_{n+1}u_{n_+1})
.\] 
\[
	2u_n - 2u_{n-1} + 3u_{n+1} - 6 u_n + 3u_{n+1} = 3u_{n+1} - 4u_n + u_{n-1} = 2h f_{n-1}
.\] 
\[
	\Rightarrow u_{n+1} = \underbrace{\frac 43}_{\alpha_0} \underbrace{-\frac 13}_{\alpha_1}u_{n-1} + \underbrace {\frac 23}_{b_{-1}}hf_{n+1}
.\]  
\begin{center}
	
\renewcommand{\arraystretch}{1.5} % Aumenta lo spazio verticale per le frazioni
\begin{tabular}{|c|c|c|c|c|c|c|c|c|}
\hline
& $P_0$ & $a_0$ & $a_1$ & $a_2$ & $a_3$ & $a_4$ & $a_5$ & $b_{-1}$ \\
\hline
BDF1 & $0$ & $1$ & $-$ & $-$ & $-$ & $-$ & $-$ & $1$ \\
\hline
BDF2 & $1$ & $\frac{4}{3}$ & $-\frac{1}{3}$ & $-$ & $-$ & $-$ & $-$ & $\frac{2}{3}$ \\
\hline
BDF3 & $2$ & $\frac{18}{11}$ & $-\frac{9}{11}$ & $\frac{2}{11}$ & $-$ & $-$ & $-$ & $\frac{6}{11}$ \\
\hline
BDF4 & $3$ & $\frac{48}{25}$ & $-\frac{36}{25}$ & $\frac{16}{25}$ & $-\frac{3}{25}$ & $-$ & $-$ & $\frac{12}{25}$ \\
\hline
BDF5 & $4$ & $\frac{300}{137}$ & $-\frac{300}{137}$ & $\frac{200}{137}$ & $-\frac{75}{137}$ & $\frac{12}{137}$ & $-$ & $\frac{60}{137}$ \\
\hline
BDF6 & $5$ & $\frac{360}{147}$ & $-\frac{450}{147}$ & $\frac{400}{147}$ & $-\frac{225}{147}$ & $\frac{72}{147}$ & $-\frac{10}{147}$ & $\frac{60}{147}$ \\
\hline
\end{tabular}
\end{center}
\end{table}
\begin{teo}
	I metodi BDF sono consistenti, con ordine pari al numero di passi
\end{teo}
\begin{teo}
	I metodi BDF sono $\theta$-stabili solo per $p\leq 5$
\end{teo}
\textbf{Osservazione}\\
I BDF sono sempre impliciti quindi vanno risolti con metodi per la ricerca di zeri (punto fisso o newton)\\
\begin{teo}
	I metodi BDF sono $\theta$-stabili solo per $p \leq 5$\\
	La loro regione di stabilità assoluta include tutto l'asse negativo ($A(\theta)$-stabili)
\end{teo}
\textbf{Esempio Stiff}:\\
Soluzione: \begin{cases}
	$y_1 = e^{-\lambda_1 t}\\
	y_2 = e^{-\lambda_2 t}$
\end{cases} $\rightarrow \begin{cases}
	\dot y_1 = -\lambda_1 y_1\\
	\dot y_2 = - \lambda_2 y_2\\
	y_1(0)\\
	y_2(0)
\end{cases}$\\
Se $0 < \lambda_1 << \lambda_2$
Se uso FE: $|1 + \lambda_1 h| < 1, \ \ |1 + \lambda_2h| < 1\leadsto $ questa è più restrittiva\\
\textbf{Zero-Stabilità}\\
$p =1 \rightarrow BDF 2: \rho(r) = r^{p+1} - \sum^{p}_{j=0}a_jr^{p-j} = r^2 - \frac 43 r + \frac 13 = (r-1)(r-\frac 13) \leadsto r_1 = 1, \ r_2= \frac 13$\\
\begin{defi}[Boundary Lotus]
	Chiamiamo la regione di stabilità assoluta per un LMM Boundary Lotus
\end{defi}
\[
	P(r) = \rho(r) - z\sigma(r) = r^{p+1} - \sum^{p}_{j=0}a_jr^{p-j} - z \sum^{p}_{j=-1}b_jr^{p-j}
.\] 
$z\in A \Leftrightarrow |r_j(z)| < 1$ per $j = 0,\ldots, p$ con  $P(r_j(z)) = 0$\\
Possiamo esplicitare  $P(r)$ rispetto a $z$ ovvero se $r$ è radice $P(r) = 0$
 \[
	 P(r) = 0 \Leftrightarrow\rho(r) - z\sigma(r) = 0 \Leftrightarrow z = \frac{\rho(r)}{\sigma(r)}= \frac{r^{p+1} - \sum^{p}_{j=0}a_jr^{p-j}}{ \sum^{p}_{j=-1}b_jr^{p-j}} \ \ \ \circledast
\] 
Se $z \in \partial A \Rightarrow  \exists \bar j \ t.c.\ |r_{\bar j}(z)| = 1 \Leftrightarrow r_{\bar j}(z) = e^{i\bar \theta}$ per qualche $\bar \theta\in[0,2\pi]$\\
Se sostituiamo $r_{\bar j}(z)$ su $\circledast$
 \[
	 z = \frac{e^{i(p+1)\bar\theta} - \sum^{p}_{j=0}a_je^{i(p-j)\bar\theta}}{ \sum^{p}_{j=-1}b_je^{i(p-j)\bar\theta}} = z(\bar\theta)
.\] 
Al variare $\bar \theta$ $z(\bar\theta)$ descrive una curva nel piano complesso
 \[
	 \Rightarrow \partial A\subseteq \{z(\bar \theta) \ \ t.c.\ \ \bar \theta\in[0,2\pi]\}
.\] 

\end{document}
